% sintaxismark-doc-en.tex
% English documentation for sintaxismark.sty v0.7.2 experimental monolithic
% Compile with LuaLaTeX when using sintaxisgrid.

\documentclass[11pt,a4paper]{article}

\usepackage{fontspec}
\usepackage[english]{babel}
\usepackage{geometry}
\geometry{margin=2.5cm}
\usepackage{booktabs}
\usepackage{longtable}
\usepackage{xcolor}
\usepackage{listings}
\usepackage{hyperref}
\usepackage{parskip}
\usepackage{microtype}
\usepackage{gb4e}\noautomath

\usepackage[docencia]{sintaxismark}

\lstset{
  language=[LaTeX]TeX,
  basicstyle=\ttfamily\small,
  columns=fullflexible,
  breaklines=true,
  frame=single,
  backgroundcolor=\color{gray!8},
  keywordstyle=\color{blue!70!black},
  commentstyle=\color{gray!70!black}
}

\title{\texttt{sintaxismark}: documentation}
\author{Pablo Damián Zdrojewski}
\date{Experimental monolithic version 0.7.2 --- 2026-06-29}

\begin{document}
\maketitle
\tableofcontents

\section{Overview}

\texttt{sintaxismark} is a LaTeX package for marking syntactic constituents, grammatical functions, and heads directly over running text. It is intended for grammar teaching materials, exercises, slides, and compact linguistic analyses.

Experimental version 0.7.2 uses one main file, \texttt{sintaxismark.sty}. This file contains both the classic mode and the experimental \texttt{grid} mode. To use the monolithic version, copy \texttt{sintaxismark.sty} to the same directory as your document.

\section{Two modes}

\begin{enumerate}
  \item \textbf{Classic mode}: uses TikZ overlays with commands such as \verb|\spanline|, \verb|\openbox|, \verb|\tagbelow|, \verb|\SES|, \verb|\PVS|, \verb|\OD|, \verb|\MI|, etc. It is suitable for short examples or for carefully controlled manual markup.
  \item \textbf{Grid mode}: uses the \verb|sintaxisgrid| environment. It reads a nested syntactic structure, extracts tokens, computes line breaks, and draws the whole analysis on a grid. It is intended for long or multiline examples because it avoids diagonal lines when a mark crosses a line break.
\end{enumerate}

Classic mode is preserved for compatibility. Grid mode is experimental and requires LuaLaTeX.

\section{Loading the package}

\begin{lstlisting}
\usepackage{sintaxismark}
\end{lstlisting}

With \texttt{gb4e}, load \texttt{gb4e} first:

\begin{lstlisting}
\usepackage{gb4e}
\noautomath
\usepackage{sintaxismark}
\end{lstlisting}

A classic visual preset can also be selected:

\begin{lstlisting}
\usepackage[docencia]{sintaxismark}
\usepackage[impresion]{sintaxismark}
\usepackage[diapositivas]{sintaxismark}
\usepackage[publicacion]{sintaxismark}
\end{lstlisting}

The presets adjust classic mode. In this experimental version, grid mode uses internal spacing values that were tuned to approximate the classic settings, especially \verb|liney|, \verb|boxd|, \verb|boxsep|, \verb|tagy|, and \verb|nestedgap|.

\section{Minimal classic example}

\begin{lstlisting}
\documentclass{article}
\usepackage{gb4e}\noautomath
\usepackage[docencia]{sintaxismark}

\begin{document}

\begin{exe}
\ex \SES{Juan} \PVS{\NV{compró} \OD{manzanas}}.
\end{exe}

\end{document}
\end{lstlisting}

Because classic mode uses TikZ \texttt{remember picture}, it is usually best to compile twice.

\section{Main shortcuts}

\begin{longtable}{ll}
\toprule
Command & Visual function \\
\midrule
\verb|\SES{...}| & simple explicit subject, upper span \\
\verb|\ST{...}| & null subject, upper span \\
\verb|\SEC{...}| & compound explicit subject, upper span \\
\verb|\PVS{...}| & simple verbal predicate, upper span \\
\verb|\PVC{...}| & compound verbal predicate, upper span \\
\verb|\OS{...}| & sentence, upper span \\
\verb|\OD{...}| & direct object, lower label with underline \\
\verb|\OD*{...}| & direct object, lower open box \\
\verb|\OI{...}| & indirect object, lower label with underline \\
\verb|\OI*{...}| & indirect object, lower open box \\
\verb|\MD{...}| & direct modifier, lower label \\
\verb|\MI{...}| & indirect modifier, lower label \\
\verb|\MI*{...}| & indirect modifier, lower open box \\
\verb|\N{...}| & head, lower label \\
\verb|\NN{...}| & nominal head, lower label with underline \\
\verb|\NV{...}| & verbal head, lower label with underline \\
\bottomrule
\end{longtable}

\section{Experimental grid mode}

\subsection{Basic example}

\begin{lstlisting}
\begin{sintaxisgrid}[width=10cm,valign=base]
\SA{
  \SES{Juan}
  \PVS{
    saw
    \OD*{
      \MD{a}
      \N{tiger}
    }
  }
}
\end{sintaxisgrid}
\end{lstlisting}

Indentation has no structural value. Structure is expressed by commands and braces. The same example can be written on one line:

\begin{lstlisting}
\begin{sintaxisgrid}[width=10cm]
\SA{\SES{Juan}\PVS{saw \OD*{\MD{a}\N{tiger}}}}
\end{sintaxisgrid}
\end{lstlisting}

\subsection{Use with \texttt{gb4e}}

\begin{lstlisting}
\begin{exe}
\ex
\begin{sintaxisgrid}[width=.92\linewidth,valign=base,yoffset=0pt]
\SA{\SES{Juan}\PVS{saw \OD*{\MD{a}\N{tiger}}}}
\end{sintaxisgrid}
\end{exe}
\end{lstlisting}

\verb|valign=base| aligns the environment with the example baseline. If the \texttt{gb4e} number is slightly too high or too low, use \verb|yoffset|:

\begin{lstlisting}
\begin{sintaxisgrid}[width=.92\linewidth,valign=base,yoffset=0.2ex]
...
\end{sintaxisgrid}
\end{lstlisting}

\subsection{Environment options}

In the current monolithic version, the stable environment options are:

\begin{longtable}{ll}
\toprule
Option & Function \\
\midrule
\verb|width=<dimension>| & available width for computing line breaks \\
\verb|valign=base| & baseline alignment, recommended with \texttt{gb4e} \\
\verb|valign=t| & top alignment \\
\verb|valign=c| & centered alignment \\
\verb|valign=b| & bottom alignment \\
\verb|yoffset=<dimension>| & fine manual vertical adjustment \\
\bottomrule
\end{longtable}

\subsection{Visual spacing}

The grid renderer uses internal values tuned to resemble classic mode. In particular:

\begin{itemize}
  \item upper spans are placed with spacing comparable to \verb|liney| and \verb|labsep|;
  \item lower boxes are placed with depth comparable to \verb|boxd| and \verb|boxsep|;
  \item simple lower tags use spacing comparable to \verb|tagy|;
  \item nested levels attempt to reproduce the effect of \verb|nestedgap| and the automatic vertical reservation used in classic mode.
\end{itemize}

Since grid mode computes one complete box, it does not depend on TikZ \texttt{remember picture} and does not need two compilation passes for placement. Warnings such as \texttt{Missing character: There is no U+000A in font nullfont} may appear in this experimental version and do not prevent PDF generation.

\section{Included demo}

The file \texttt{demo-monolithic.tex} tests the monolithic version with \texttt{gb4e}, classic mode, and grid mode. Compile with:

\begin{lstlisting}
lualatex demo-monolithic.tex
\end{lstlisting}

To check that the single file is self-contained, copy only \texttt{sintaxismark.sty} and \texttt{demo-monolithic.tex} to a temporary directory and compile there.

\section{Version notes}

\begin{itemize}
  \item \textbf{0.7.2 experimental}: monolithic version. \texttt{sintaxismark.sty} contains classic mode and grid mode with embedded Lua.
  \item \textbf{0.7.1 experimental}: fixes body capture in \verb|sintaxisgrid|.
  \item \textbf{0.7.0 experimental}: introduces \verb|sintaxisgrid|.
\end{itemize}

\end{document}
